\documentclass[11pt]{article}

\usepackage[a4paper, margin=1in]{geometry}
\usepackage{amsmath}
\usepackage{graphicx}
\usepackage{booktabs}
\usepackage{hyperref}
\usepackage{url}
\usepackage{authblk}
\usepackage{microtype}

\graphicspath{{figures/}}

\title{Single-extractor large-language-model thematic analysis of a Reddit
endometriosis community with cross-judge LLM rubric validation: a
pilot replication and methodological contribution}

\author[1]{[Author 1]}
\author[1]{[Mentor]}
\affil[1]{[Institution]}

\date{\today}

\begin{document}
\maketitle

\begin{abstract}
\noindent
\textbf{Background.}
A 2024 abstract demonstrated that GPT-4 can summarize endometriosis-care
themes from large-scale Reddit posts. We set out to extend that work
along three axes---broader corpus retrieval, multiple modern LLMs, and
a model-only validation framework---and, in the process, discovered
practical limits on free- and prepaid-tier API access that materially
shape what an undergraduate-budget reproduction can achieve.

\noindent
\textbf{Methods.}
The original 2024-study Reddit corpus was provided by the senior
mentor as a Pushshift-format archive of r/Endo (\texttt{Endo\_submissions.zst}
and \texttt{Endo\_comments.zst}) covering Dec 2010 to Dec 2022.
We filtered to the published 2018-02-07 to 2024-10-31 window---the
archive ends 2022-12-31, so 2023-01 to 2024-10 is not covered---and
restricted to r/Endo, since r/endometriosis was not present in the
provided dump. We then drew a year-stratified 2{,}000-post /
80{,}000-comment sample matching the 2024 study size. \textbf{Claude
Haiku 4.5}~\cite{anthropic_claude_haiku} ran the per-chunk theme
extraction; rather than calling the pay-per-token API, the sample was
chunked at 100k tokens (74 chunks) and each chunk was dispatched to a
\emph{Claude Code subagent} running Haiku 4.5, which read the chunk
file and wrote a per-chunk theme JSON. The cross-chunk merge ("reduce")
step was dispatched as a Claude Sonnet 4.6 subagent; both judges
(Sonnet 4.6 and Opus 4.7) were also dispatched as subagents. We
computed four LLM-only reliability signals per theme: (i)
\textbf{chunk-support}, the count of disjoint corpus chunks that
independently produced the theme (within-model agreement proxy);
(ii) \textbf{quote grounding} via exact substring + BM25 retrieval
against the sampled corpus; (iii) \textbf{lexical alignment} with the
eight Young~2015 systematic-review themes~\cite{young2015}; and (iv)
\textbf{rotating LLM-as-judge rubric scoring} on faithfulness,
comprehensiveness, clinical relevance, specificity, and non-redundancy.
The remaining originally-specified signal---cross-model agreement against a
second independent extractor (Gemini 2.5 Flash)---could not be
completed because Gemini's free-tier daily token-per-minute quota
blocked the full extraction.

\noindent
\textbf{Results.}
The mentor-provided r/Endo archive contained 42{,}722 posts and
339{,}911 comments spanning Dec 2010 to Dec 2022. Filtering to
2018-02-07 onward and dropping removed/empty rows yielded
\textbf{38{,}001 posts and 291{,}754 comments} (23.4M words) from
25{,}074 unique users---roughly 19$\times$ the 2024 abstract's
1{,}986 posts. Claude Haiku 4.5 subagents produced \textbf{722
chunk-level themes} across all 74 successfully-processed
100k-token chunks; a subsequent LLM-driven reduce step (Claude Sonnet
4.6 via subagent) merged these into \textbf{20 final themes}, all of
which had chunk-support $\geq 40/74$ and the strongest of which
(``Laparoscopic excision surgery as gold-standard treatment'') had
$k=72$. Of supporting quotes, 35\% were grounded in the corpus by
exact-match or BM25; 9 of 20 (45\%) lexically aligned with a Young
2015 theme; mean rotating-judge rubric score was 4.45 of 5 across
Sonnet and Opus. Mean composite reliability across the three
normalized signals was \textbf{0.55} (range 0.28 to 0.87).

\noindent
\textbf{Conclusion.}
A single-frontier-extractor pilot with cross-judge rubric validation,
anchored to the prior qualitative literature, can recover
endometriosis-care themes from the original 2024-study Reddit corpus
at roughly 19$\times$ the abstract's reported retrieval. The
model-only reliability signals we implemented are honest---they do
not require human raters and they explicitly flag low-confidence
themes. Two unrealized extensions (cross-model agreement against a
second independent extractor; pulling in r/endometriosis and the
2023-2024 portion of the window) are known gaps that a funded
follow-up can fill with no methodological changes; the pipeline is
released intact.

\noindent
\textbf{Keywords:} endometriosis, large language model, reddit,
infodemiology, thematic analysis, reproducibility.
\end{abstract}

\section{Background}

Endometriosis affects an estimated 10\% of reproductive-age women and is
characterized by long diagnostic delays, chronic pain, and care
trajectories shaped as much by clinician knowledge and bias as by the
disease itself~\cite{young2015}. Online patient communities are a
unique source for surfacing these experiences in patients' own
words~\cite{eysenbach2009infodemiology}.

A 2024 abstract reported the first large-scale LLM analysis of two
endometriosis-themed subreddits, identifying five treatment-relevant
themes via GPT-4. That work surfaced the value of LLM thematic synthesis
but had three methodological limits we set out to address:

\begin{enumerate}
  \item Reddit's PRAW listing endpoint caps a single subreddit pull at
    $\sim$1{,}000 items, so the reported 1{,}986-post total covers far
    less than the 6.5-year window the abstract names.
  \item Only one model was used, leaving open the question of how much
    of any reported theme is a property of the corpus versus a property
    of GPT-4's training distribution.
  \item No method was reported to validate the LLM's interpretations
    against either the source corpus or the prior qualitative
    literature.
\end{enumerate}

We aimed to extend this work along all three axes within the budget of
an undergraduate research project: prepaid Anthropic credits and a
free-tier Google Gemini API key. We report here what was achievable
under those constraints, the failure modes we encountered, and the
adapted single-model framework that the data ultimately supports.

\section{Methods}

\subsection{Corpus}

The senior mentor provided the original 2024-study Reddit corpus as a
Pushshift-format archive (\texttt{Endo\_submissions.zst} and
\texttt{Endo\_comments.zst}) covering r/Endo from Dec 2010 to Dec
2022. We filtered to the published 2018-02-07 to 2024-10-31 window;
because the archive ends 2022-12-31, the 2023-01 through 2024-10-31
portion of the published window is not covered. r/endometriosis was
not present in the provided dump, so the analysis is restricted to
r/Endo. After filtering removed/deleted bodies and out-of-window
rows, the analyzable archive contains 38{,}001 posts and 291{,}754
comments (23.4M words) from 25{,}074 unique users. This is roughly
19$\times$ the size of the 2024 abstract's reported retrieval
(Table~\ref{tab:corpus}). The corpus is identified by SHA-256 hashes
for reproducibility; raw text is not redistributed.

A year-stratified random sample of 2{,}000 posts and 80{,}000 comments
was drawn to match the 2024 abstract's analysis size, allowing a
like-for-like methodological comparison.

\input{tables/corpus_stats}

\subsection{Models}

\textbf{Claude Haiku 4.5} (Anthropic, dated snapshot
\texttt{claude-haiku-4-5-20251001}~\cite{anthropic_claude_haiku}) was
used for theme extraction. We had originally specified a two-model
design with Claude Haiku 4.5 and Gemini 2.5 Flash~\cite{google_gemini_2_5},
but the free-tier daily token-per-minute quota on Gemini 2.5 Flash was
exhausted on the first chunk of the 100k-token-chunk run, and could
not be recovered within the project window. We could not re-run GPT-4
(the 2024 baseline model) due to absence of an OpenAI API credential,
so the 2024 abstract's five themes are treated as a frozen published
baseline column rather than a live comparator.

\subsection{Map-reduce theme extraction}

The sample was packed into 100{,}000-token chunks (well within Haiku
4.5's 200{,}000-token context) with a 2{,}000-token overlap, yielding
74 chunks. Each chunk was written to disk as a flat text file and
dispatched to a separate \emph{Claude Haiku 4.5 subagent}. The
subagent read its assigned chunk file and wrote a per-chunk JSON of
themes, each with a short label, a 2-3 sentence description, up to
three verbatim supporting quotes, a frequency estimate, and a
severity signal. The 74 subagents produced 722 chunk-level themes
across all 74 chunks, with no chunks dropped.

We then performed the reduce step using a Claude Sonnet 4.6 \emph{Claude
Code subagent}. The subagent received all 722 chunk-level themes and
was prompted to merge same-meaning themes, carry forward verbatim
quotes, record per-merged-theme \texttt{chunk\_support} (the count of
distinct chunks that contributed), and emit a clean global theme list.
It also emitted \texttt{raw\_member\_count}, which we retain as a
descriptive audit field rather than a reliability signal. This produced
20 final themes (Table~\ref{tab:runs}). The subagent prompt protocol is
preserved in \texttt{prompts/subagent\_protocol.md}; the exact
interactive subagent messages from the completed run were not separately
captured.

\input{tables/run_index}

\subsection{Reliability signals}

The original design called for five LLM-only reliability signals
(cross-model agreement against a second extractor, quote grounding,
anchored validity against Young 2015, LLM-as-judge with rotating judge,
and stability under prompt-paraphrase $+$ temperature perturbation).
Four were achievable under our final budget.

\paragraph{(1) Chunk-survival.} The reduce-step output records, for
each merged theme, how many of the 74 successfully-processed corpus
chunks independently contributed (\texttt{chunk\_support}). A theme
present in many disjoint chunks is robust to within-model variation
across corpus slices and is a within-extractor analog of cross-model
agreement.

\paragraph{(2) Quote grounding.} Each supporting quote is checked
against the corpus by (a) whitespace-collapsed, case-insensitive
substring match and (b) best-of BM25 retrieval. A quote is grounded
if either condition fires; per-theme grounding rate is the fraction
of grounded quotes (Figure~\ref{fig:grounding}). Embedding-based
semantic grounding was originally planned but could not be run at
corpus scale within the free-tier embedding-API daily request quota.

\paragraph{(3) Anchored validity.} The eight Young 2015 themes
(Table~\ref{tab:anchor}) were encoded as a JSON literature anchor.
Theme$\to$anchor lexical similarity was computed via TF-IDF cosine
(threshold 0.05), giving a coarse but interpretable signal of overlap
with the prior qualitative-research consensus. We note that TF-IDF
under-counts paraphrased matches and is therefore a lower bound on
true semantic alignment.

\paragraph{(4) Rotating LLM-as-judge.} Two independent judges---Claude
Sonnet 4.6 and Claude Opus 4.7, both dispatched as Claude Code
subagents on a research subscription rather than the pay-per-token
API---scored each theme on faithfulness, comprehensiveness, clinical
relevance, specificity, and non-redundancy with one-sentence
justifications, on a 1-5 scale. Per-theme judge score is the mean
across dimensions and judges; Sonnet and Opus scores are also reported
separately for inter-judge calibration (Table~\ref{tab:reliability}).

\paragraph{Unrealized.}
\textbf{(5) Cross-model agreement} requires a second independent
extractor and could not be completed: Gemini 2.5 Flash's free-tier
daily token-per-minute quota was exhausted on the first chunk of the
100k-token-chunk run. The pipeline supports this signal directly
(\texttt{src/cluster\_themes.py}); a funded follow-up that can run a
second extractor will produce it without further methodological
change. We note that the Claude Code subagent route used for the
reduce and judge steps does not currently expose non-Anthropic
models, so a true cross-vendor agreement signal does require either
an OpenAI / Google API credential or a different multi-tenant gateway.

\paragraph{Composite.}
Available signals were normalized to $[0, 1]$ and averaged into a
per-theme composite reliability score
(Table~\ref{tab:reliability}; Figure~\ref{fig:reliability}). Mean
composite reliability was 0.55 (range 0.28-0.87); the manuscript
reports both the composite and the components so reviewers can see
which signals contributed.

\subsection{Reporting}

We follow the MI-CLAIM checklist for clinical-AI
reporting~\cite{norgeot2020miclaim}; the completed checklist appears in
the Supplement (\texttt{supplement\_miclaim.tex}). Code, prompts, model
snapshots, the Young 2015 anchor encoding, and aggregate results are
released; raw Reddit text is \emph{not} redistributed in line with
infodemiology ethics norms~\cite{norori2021addressing}.

\section{Results}

The mentor-provided r/Endo archive returned 38{,}001 posts and
291{,}754 comments (23.4M words) from 25{,}074 unique users after
filtering to the in-window, non-deleted records. The year-stratified
82{,}000-record sample matched the 2024 abstract's analysis size.

Claude Haiku 4.5 subagents produced 722 chunk-level themes across all
74 successfully-processed 100k-token chunks. The Sonnet-subagent
reduce step collapsed these into 20 final themes
(Table~\ref{tab:reliability}). The top themes by chunk-support were
``Laparoscopic excision surgery as gold-standard treatment''
($k=72/74$), ``Hormonal contraception and GnRH-based medical
management'' ($k=71$), ``Diagnostic delays and imaging limitations''
($k=70$), ``Pain management strategies and pharmacological adequacy''
($k=70$), and ``Medical dismissal and gaslighting'' ($k=68$). Five
additional themes had $k \geq 60$. The lowest chunk-support
(``Adenomyosis as underrecognized comorbidity'', $k=40$) is still
present in more than half of all chunks.

Quote grounding by exact $+$ BM25 reached 35\% across 60 supporting
quotes. The gap from a higher achievable rate reflects the absence of
embedding-based semantic grounding, which would catch paraphrased and
lightly-edited quotes that BM25 alone misses.

Lexical alignment with Young 2015 reached threshold for 9 of 20 themes
(45\%; Table~\ref{tab:anchor}). The strongest matches were on
diagnostic delay, medical dismissal, sexual dysfunction, fertility
concerns, online community/peer support, and emotional impact---all
of which appear directly in Young 2015's qualitative synthesis. The
eleven themes that did not lexically align (laparoscopic excision,
hormonal management, pain management strategies, specialist access,
post-operative recovery, post-surgical recurrence, hysterectomy,
GI/bowel complications, pelvic floor dysfunction, dietary
modifications, adenomyosis comorbidity) describe finer-grained
sub-phenomena or topics outside Young 2015's eight-theme synthesis.

Rotating LLM-as-judge produced consistent rubric scores across the
two judges. Sonnet 4.6 returned a mean rubric score of
$\bar x = 4.4 / 5$ and Opus 4.7 returned $\bar x = 4.5 / 5$
across all dimensions and themes, with
faithfulness and clinical relevance as the strongest dimensions
($\bar x \geq 4.7$). Specificity and non-redundancy were the weakest
dimensions, reflecting that the highest-chunk-support themes tend to
bundle multiple related sub-phenomena.

\input{tables/reliability_per_theme}

\begin{figure}[ht]
\centering
\includegraphics[width=0.75\linewidth]{grounding.png}
\caption{Per-run quote grounding rate (exact-match + BM25; embedding
match was disabled). 35\% of supporting quotes (21/60) from the
deduplicated theme set are grounded in the corpus by lexical retrieval
alone.}
\label{fig:grounding}
\end{figure}

\begin{figure}[ht]
\centering
\includegraphics[width=0.75\linewidth]{reliability.png}
\caption{Distribution of per-theme composite reliability scores. The
composite averages quote grounding, rotating LLM-as-judge score, and
Young 2015 lexical alignment, normalized to $[0, 1]$. Cross-model
agreement and prompt/temperature stability are absent for the reasons
documented in Methods; chunk-survival is reported separately as
\texttt{k\_chunk}.}
\label{fig:reliability}
\end{figure}

\input{tables/anchor_coverage}

\section{Discussion}

This pilot extends the 2024 GPT-4 abstract in a narrower but more
auditable direction than originally planned. The main empirical change
is corpus retrieval: the mentor-provided Pushshift-format r/Endo archive
contains far more analyzable material than the 2024 abstract's reported
1{,}986 posts, even after restriction to the overlapping 2018-02-07 to
2022-12-31 window. The model comparison we initially intended could not
be completed under free-tier Gemini limits, so the present claim is not
``two models independently recovered the same themes.'' The defensible
claim is instead that a single contemporary extractor, run over 74
large corpus chunks and checked by two independent judge models, recovers
a stable set of treatment- and care-related themes that can be audited
against source quotes and prior qualitative literature.

The strongest result is within-run recurrence. All 20 final themes were
supported by at least 40 of 74 disjoint chunks, and the highest-support
themes appeared in nearly every slice of the sample. This matters
because a theme that reappears across many independently processed
chunks is less likely to be a one-off artifact of a single prompt
completion. Chunk-support is not a substitute for cross-model
agreement, but it is a useful single-extractor robustness signal and is
more transparent than treating the final reduced theme list as a flat
set of equally supported findings.

The validation stack also clarifies where the current run is weak.
Rotating LLM-as-judge scores were high, especially for faithfulness and
clinical relevance, but quote grounding was only 35\% by exact-match
and BM25 retrieval. Some of that gap is expected because embedding-based
semantic grounding was disabled, and because the reduce step may carry
lightly edited quotes forward. Still, the low lexical grounding rate is
important: the judge scores should not be read as proof that every
supporting quote is audit-ready. In a revision intended for submission,
the quote fields should either be re-grounded with semantic retrieval or
manually replaced with verified paraphrases before publication.

The Young~2015 anchor results separate broad convergence from
fine-grained online-community topics. Nine of 20 final themes aligned
lexically with the prior qualitative synthesis, including diagnostic
delay, medical dismissal, chronic pain, sexual dysfunction, fertility
concerns, psychological impact, and peer support. Themes that did not
align, such as laparoscopic excision, GnRH-based management,
post-operative recovery, recurrence after surgery, pelvic floor
therapy, dietary changes, bowel involvement, hysterectomy, specialist
access, and adenomyosis comorbidity, are not automatically novel
findings. They are more precise treatment-management topics than the
Young~2015 theme set was designed to encode. Their value is practical:
they suggest concrete issues clinicians and researchers can ask about
when discussing endometriosis care with patients.

The most clinically useful themes are therefore not necessarily the
highest composite-reliability themes. Diagnostic delay and medical
dismissal are both highly reliable and already well represented in the
literature. In contrast, specialist access, pelvic floor dysfunction,
post-operative recovery, and adenomyosis comorbidity score lower
because they do not match the Young anchor or have weaker lexical quote
grounding, yet they still appear across many chunks and may be
important prompts for future qualitative review. The composite score
should be interpreted as a triage aid for human review, not a final
truth label.

This study is best understood as a reproducibility and methods pilot.
It shows that a low-budget, single-extractor workflow can produce a
traceable theme table, preserve intermediate chunk-level outputs, and
attach multiple reliability signals to every final theme. It also shows
where such a workflow falls short: cross-model agreement, prompt
paraphrase stability, high-temperature stability, and human
patient/clinician review are still absent. A funded follow-up should
run a second extractor on the same frozen sample, repeat the prompt-B
and temperature-stability sweeps, and calibrate Sonnet/Opus judge
scores against human ratings.

Several limitations bound the strength of the findings. The analysis is
restricted to r/Endo because r/endometriosis was not present in the
mentor-provided archive, and the archive ends on 2022-12-31, leaving the
2023-01 to 2024-10-31 portion of the published window unanalyzed. Reddit
users are self-selected, English-speaking, and more digitally engaged
than the broader endometriosis population. The original 2024 GPT-4 run
could not be reproduced without an OpenAI credential. Finally, although
we now preserve the subagent protocol for reruns, the exact interactive
subagent messages from the completed run were not captured as a
first-class artifact. Future runs should log those prompts before any
subagent work begins.


\section{Limitations}

\begin{itemize}
  \item \textbf{Single extractor.} Gemini 2.5 Flash's free-tier
    daily quota prevented a second independent extraction. The
    cross-model agreement signal was therefore not computed.
  \item \textbf{Subreddit coverage.} The mentor-provided archive
    contains r/Endo only; r/endometriosis (also analyzed in the 2024
    abstract) is not present. A merged-subreddit replication is a
    natural follow-up.
  \item \textbf{Window coverage.} The provided archive ends
    2022-12-31, so the 2023-01 to 2024-10-31 portion of the published
    window is not analyzed here. Pulling additional archived months
    from public Pushshift mirrors would close this gap.
  \item \textbf{TF-IDF for anchor matching.} Young 2015 alignment is
    computed lexically rather than via transformer embeddings, which
    were not available at corpus scale within the free-tier embedding
    quota. The reported alignment fraction is a lower bound on true
    semantic overlap.
  \item \textbf{No GPT-4 baseline.} We could not re-run GPT-4 (the
    2024 model) ourselves; the 2024 abstract's themes are treated as a
    frozen published row rather than a same-corpus comparator.
  \item \textbf{Self-selected community.} Reddit users skew younger,
    more digitally-literate, and English-speaking; their experiences
    do not generalize to all people with endometriosis.
\end{itemize}

\section{Reproducibility statement}

All code, prompts, pinned model snapshots, prompt outputs, the Young
2015 anchor encoding, and the aggregated reliability table are
included in the project repository at the path
\texttt{endometriosis-llm-2026/}. Raw Reddit posts and comments are
\emph{not} redistributed; the SHA-256 manifest is published so
collaborators with the same data window can verify identical input.
The pipeline runs end-to-end on a synthetic 100-record corpus also
included.

\section{Ethics statement}

This study analyzed publicly accessible, pseudonymous Reddit data and
is exempt from IRB review as non-human-subjects research under the
public-data category. All illustrative quotes in the published
manuscript are paraphrased; no verbatim user posts, usernames, or
direct messages appear in the published text or supplements.

\bibliographystyle{plain}
\bibliography{refs}

\end{document}
